\documentclass[a4paper,12pt]{article}
\usepackage[T2A]{fontenc}
\usepackage[utf8]{inputenc}
\usepackage[russian]{babel}
\usepackage{amsmath,amssymb,mathtools}
\usepackage{helvet}
\renewcommand{\familydefault}{\sfdefault}
\usepackage{icomma}
\usepackage{geometry}
\geometry{left=22mm,right=22mm,top=22mm,bottom=22mm}
\usepackage{microtype}
\usepackage{tikz}
\usetikzlibrary{arrows.meta,positioning,shapes.geometric,calc}
\usepackage{tabularx,booktabs,longtable,array}
\usepackage{enumitem}
\usepackage{hyperref}
\usepackage{parskip}
\usepackage{fancyhdr}
\usepackage{setspace}
\usepackage{xcolor}
\usepackage{needspace}
\definecolor{accent}{HTML}{008080}
\definecolor{darkaccent}{HTML}{005A5A}
\definecolor{textgray}{HTML}{30373B}
\definecolor{muted}{HTML}{667477}
\definecolor{linegray}{HTML}{CFDDDC}
\hypersetup{hidelinks,pdftitle={Дерево вероятностей — чек-лист для Марка Гукина},pdfauthor={Лёвин Артём Александрович}}
\onehalfspacing
\setlength{\parindent}{0pt}
\setlength{\parskip}{6pt}
\setlist[itemize]{itemsep=2pt,topsep=4pt,leftmargin=19pt}
\setlist[enumerate]{itemsep=2pt,topsep=5pt,leftmargin=19pt}
\pagestyle{fancy}
\fancyhf{}
\fancyhead[L]{\small\textcolor{muted}{Дерево вероятностей}}
\fancyhead[R]{\small\textcolor{muted}{Марк Гукин}}
\fancyfoot[C]{\small\textcolor{muted}{\thepage}}
\renewcommand{\headrulewidth}{0.25pt}
\renewcommand{\footrulewidth}{0pt}
\setlength{\headheight}{15pt}
\newcommand{\sect}[1]{\needspace{5\baselineskip}\vspace{8pt}{\Large\color{darkaccent}\bfseries #1}\par\vspace{2pt}\noindent\textcolor{linegray}{\rule{\linewidth}{0.65pt}}\vspace{3pt}}
\newcommand{\subsect}[1]{\needspace{3\baselineskip}{\large\color{darkaccent}\bfseries #1}\par\vspace{0pt}}
\newcommand{\key}[1]{\textcolor{darkaccent}{\bfseries #1}}
\newcommand{\prob}[1]{\mathrm{P}(#1)}
\newcommand{\exercise}[2]{\needspace{7\baselineskip}\vspace{7pt}\noindent{\large\color{darkaccent}\bfseries #1.}\hspace{6pt}\textbf{#2}\par\vspace{1pt}}
\newcommand{\checkedline}[1]{\noindent\(\square\)\quad #1\par}
\begin{document}
\thispagestyle{empty}
\begin{center}
\vspace*{22mm}
{\Large\color{darkaccent} МАТЕМАТИКА}\par
\vspace{17mm}
{\LARGE\color{textgray} Чек-лист}\par
\vspace{15mm}
{\huge\color{accent}\bfseries Дерево вероятностей}\par
\vspace{4mm}
{\Large\color{accent} Пути событий и полная вероятность}\par
\vspace{18mm}
{\large ЕГЭ · профильная математика}\par
\vfill
{\Large Марк Гукин}\par
\vspace{11mm}
{\large 22 ноября 2026 г.}\par
\vspace{15mm}
{\normalsize Лёвин Артём Александрович}\par
{\normalsize эксклюзивно для Марка Гукина}\par
\vspace{19mm}
\begin{tikzpicture}[x=1cm,y=1cm]
\draw[accent,line width=0.75pt] (0,0) .. controls (1.3,0.5) and (2.0,0.2) .. (3.0,0)
 .. controls (4,-0.2) and (4.8,-0.25) .. (5.9,0.07)
 .. controls (7,0.4) and (8.2,-0.25) .. (9.7,0);
\fill[accent] (5.9,0.07) circle(1.3pt);
\end{tikzpicture}
\end{center}
\clearpage

\sect{1. Самопроверка: что нужно уметь}
После повторения материала проверьте себя по каждому пункту. Поставьте отметку, если можете объяснить правило и применить его в новой задаче.
\checkedline{Различать отдельное событие, последовательность событий и несколько возможных сценариев.}
\checkedline{Строить дерево в хронологическом порядке: от первого выбора к последующим результатам.}
\checkedline{Подписывать \textbf{события у узлов}, а \textbf{вероятности на ветвях}.}
\checkedline{Находить вероятность \textbf{одного пути} умножением вероятностей его ветвей.}
\checkedline{Складывать вероятности \textbf{несовместных подходящих путей}.}
\checkedline{Использовать противоположное событие, когда прямой перебор слишком громоздок.}
\checkedline{Составлять уравнение с неизвестной долей \(x\) и проверять \(0\le x\le1\).}
\checkedline{Точно умножать десятичные дроби и пользоваться степенями десяти.}

\sect{2. Правила, на которых держится дерево}
\subsect{Один путь: умножение}
Если последовательно должны произойти события \(A\) и \(B\), то
\[
 \prob{A\cap B}=\prob{A}\cdot\prob{B\mid A}.
\]
Здесь \(\prob{B\mid A}\) означает вероятность второго события при условии, что первое уже произошло. \key{На дереве перемножаются числа, написанные на ветвях одного пути.}

При независимых событиях вероятность второго события не меняется после первого, поэтому \(\prob{B\mid A}=\prob{B}\).

\subsect{Несколько подходящих путей: сложение}
Если искомый результат получается по \textbf{разным несовместным путям} \(C_1,C_2,\ldots\), то
\[
 \prob{C_1\cup C_2}=\prob{C_1}+\prob{C_2},\qquad C_1\cap C_2=\varnothing.
\]
\key{Рабочее правило:} умножайте вдоль ветвей; складывайте готовые вероятности разных подходящих путей.

\subsect{Вероятность противоположного события}
Если легче найти вероятность события «\(A\) не произошло», используйте
\[
 \prob{A}=1-\prob{\overline A}.
\]
На каждом узле сумма вероятностей \textbf{всех исходящих ветвей} равна единице. Поэтому при двух исходах с вероятностями \(p\) и \(1-p\) можно быстро восстановить неизвестное число.

\sect{3. Короткие примеры из методов занятия}
\subsect{Две партии в шахматы}
Игрок поочерёдно играет белыми и чёрными фигурами. Вероятность победы белыми равна \(0{,}6\), чёрными — \(0{,}4\). Цвет первой партии определяется равновероятно; результаты партий при фиксированных цветах независимы. Для двух побед подходят \textit{две} последовательности цветов:
\[
 \frac12\cdot0{,}6\cdot0{,}4+\frac12\cdot0{,}4\cdot0{,}6
 =\boxed{0{,}24}.
\]
Оба случая уже учтены. Ещё один множитель \(\tfrac12\) добавлять нельзя.

\subsect{Четыре выстрела с заданными исходами}
Вероятность попадания первым выстрелом равна \(0{,}7\), каждым следующим — \(0{,}9\). Выстрелы независимы. Если требуется последовательность \(\text{Н, Н, П, П}\), получаем
\[
 (1-0{,}7)(1-0{,}9)\cdot0{,}9\cdot0{,}9
 =0{,}3\cdot0{,}1\cdot0{,}9^2=\boxed{0{,}0243}.
\]
Когда указан \textbf{единственный порядок} исходов, развёрнутое дерево часто избыточно.

\subsect{Чётность произведения цифр}
Три последние цифры телефонного номера выбираются независимо и равновероятно из \(0,1,\ldots,9\). Произведение нечётно \textbf{только тогда}, когда все три цифры нечётные. Для каждой позиции есть пять нечётных цифр из десяти:
\[
 \prob{\text{нечётно}}=\left(\frac5{10}\right)^3=0{,}125,\qquad
 \prob{\text{чётно}}=1-0{,}125=\boxed{0{,}875}.
\]
Чётное произведение возможно при одной, двух или трёх чётных цифрах, поэтому метод дополнения здесь проще.

\clearpage
\sect{4. Как выглядит корректно подписанное дерево}
\textbf{Задача.} В магазин поступают лампы с двух заводов: \(40\%\) с первого и \(60\%\) со второго. На первом заводе брак составляет \(2\%\), на втором — \(3\%\). Найдите вероятность того, что выбранная лампа исправна.

\begin{center}
\begin{tikzpicture}[
 >=Latex,thick, node distance=18mm and 36mm,
 every node/.style={font=\small},
 event/.style={inner sep=3pt,align=center},
 edge from parent/.style={draw=accent,->,line width=0.9pt},
 level 1/.style={sibling distance=76mm,level distance=22mm},
 level 2/.style={sibling distance=37mm,level distance=28mm}]
\node[event] {Случайная лампа}
 child {node[event] {Завод I}
   child {node[event] {Исправна} edge from parent node[left,font=\small] {\(0{,}98\)}}
   child {node[event] {Брак} edge from parent node[right,font=\small] {\(0{,}02\)}}
   edge from parent node[left,font=\small] {\(0{,}4\)}
 }
 child {node[event] {Завод II}
   child {node[event] {Исправна} edge from parent node[left,font=\small] {\(0{,}97\)}}
   child {node[event] {Брак} edge from parent node[right,font=\small] {\(0{,}03\)}}
   edge from parent node[right,font=\small] {\(0{,}6\)}
 };
\end{tikzpicture}
\end{center}

Подходят два пути: «завод I \textbf{и} исправна» \textbf{или} «завод II \textbf{и} исправна». Поэтому
\[
 \prob{\text{исправна}}=0{,}4\cdot0{,}98+0{,}6\cdot0{,}97
 =0{,}392+0{,}582=\boxed{0{,}974}.
\]
\key{Проверьте дерево:} возле каждого узла написан результат развилки, возле каждой стрелки — вероятность перехода. При чтении слева направо (или сверху вниз) порядок событий не меняется.

\sect{5. Обратная задача: находим неизвестную долю}
Предположим, что доля ламп завода I неизвестна и обозначена \(x\). Тогда доля второго завода равна \(1-x\). Известны доли брака \(0{,}02\) и \(0{,}03\), а общая вероятность купить бракованную лампу составляет \(0{,}024\). Записываем сумму путей, заканчивающихся браком:
\[
 \underbrace{0{,}02x}_{\text{завод I и брак}}+
 \underbrace{0{,}03(1-x)}_{\text{завод II и брак}}=0{,}024.
\]
\[
 0{,}02x+0{,}03-0{,}03x=0{,}024
 \quad\Longrightarrow\quad -0{,}01x=-0{,}006
 \quad\Longrightarrow\quad \boxed{x=0{,}6}.
\]
\textbf{Проверка:} \(0{,}6\cdot0{,}02+0{,}4\cdot0{,}03=0{,}024\). Доля \(0{,}6\) соответствует \(60\%\) и лежит между \(0\) и \(1\).

\sect{6. Контроль качества: событие и решение контролёра}
Исправная батарейка может быть ошибочно забракована, а неисправная — выявлена правильно. Если бракованных батареек \(2\%\), исправных \(98\%\), контролёр бракует исправную с вероятностью \(0{,}01\), а неисправную — с вероятностью \(0{,}99\), то
\[
 \prob{\text{забракована}}
 =0{,}98\cdot0{,}01+0{,}02\cdot0{,}99
 =0{,}0098+0{,}0198=\boxed{0{,}0296}.
\]
Это два разных способа попасть в категорию «забракована»: ошибка контроля \textbf{или} верное обнаружение брака. Не путайте \textbf{состояние изделия} с \textbf{решением системы проверки}.

\sect{7. Вычисления с десятичными дробями}
\subsect{Через степень десяти}
\[
 0{,}7\cdot0{,}015=(7\cdot10^{-1})(15\cdot10^{-3})
 =105\cdot10^{-4}=\boxed{0{,}0105}.
\]
\subsect{Сложение после вынесения общего множителя}
\[
 \begin{aligned}
 0{,}98\cdot0{,}01+0{,}02\cdot0{,}99
 &=(98+198)\cdot10^{-4}\\
 &=296\cdot10^{-4}=0{,}0296.
 \end{aligned}
\]
\textbf{Альтернатива:} при умножении столбиком сначала перемножьте числа без запятых, затем отсчитайте справа число десятичных знаков, равное их сумме в обоих множителях.

\sect{8. Типичные ошибки — проверяйте их первыми}
\begin{itemize}
\item \textbf{«И» и «или» перепутаны.} Умножение выполняется \emph{внутри пути}; сложение — \emph{между несовместными путями}.
\item \textbf{Сложены пересекающиеся случаи.} Формула простого сложения применима только к несовместным событиям (как разным исходам одного дерева).
\item \textbf{Вместо \(1-x\) поставлено \(1+x\).} Сумма долей всех поставщиков обязательно равна единице.
\item \textbf{Неправильная подпись узлов.} Подписывайте событие у узла и вероятность возле ветви; проверяйте сумму исходящих чисел.
\item \textbf{Перепутан брак с результатом контроля.} Вероятность «изделие неисправно» отличается от «изделие забраковано».
\item \textbf{Пропущен ноль в десятичной дроби.} Например, \(105\cdot10^{-4}=0{,}0105\), а \(0{,}875\) — точный результат, а не \(0{,}87\).
\end{itemize}

\clearpage
\thispagestyle{fancy}
\begin{center}{\Large\color{darkaccent}\bfseries Алгоритм: как решить задачу по дереву вероятностей}\end{center}
\vspace{4mm}
\begin{center}
\begin{tikzpicture}[
 >=Latex, node distance=7mm and 11mm,
 startstop/.style={ellipse,draw=accent,line width=0.8pt,minimum width=3.3cm,minimum height=8mm,align=center,font=\small},
 action/.style={rectangle,rounded corners=2pt,draw=accent,line width=0.75pt,minimum width=10.7cm,text width=10.0cm,minimum height=12mm,align=center,font=\small,inner sep=5pt},
 decision/.style={diamond,aspect=3.2,draw=accent,line width=0.8pt,text width=5.0cm,minimum height=13mm,align=center,font=\small,inner sep=1pt},
 branch/.style={rectangle,rounded corners=2pt,draw=accent,line width=0.75pt,text width=3.6cm,minimum height=12mm,align=center,font=\small,inner sep=4pt},
 arrow/.style={-{Latex[length=2.5mm]},line width=0.8pt,draw=darkaccent}]
\node[startstop] (s) {Начало};
\node[action,below=of s] (read) {Прочитайте условие и определите, в каком порядке происходят события};
\node[action,below=of read] (tree) {Постройте дерево: события — в узлах, вероятности — на ветвях};
\node[action,below=of tree] (paths) {Выделите все пути к нужному результату};
\node[action,below=of paths] (mult) {Для каждого пути перемножьте вероятности ветвей};
\node[action,below=of mult] (add) {Сложите вероятности всех подходящих несовместных путей};
\node[decision,below=9mm of add] (unk) {Есть неизвестная доля \(x\)?};
\node[branch,below left=11mm and 1mm of unk] (yes) {Найдите \(x\) из уравнения с \(x\) и \(1-x\)};
\node[branch,below right=11mm and 1mm of unk] (no) {Запишите полученную численную вероятность};
\node[action,below=15mm of unk,yshift=-18mm] (check) {Проверьте: сумма на развилках равна 1; ответ лежит от 0 до 1};
\node[startstop,below=of check] (end) {Ответ};
\draw[arrow] (s)--(read);
\draw[arrow] (read)--(tree);
\draw[arrow] (tree)--(paths);
\draw[arrow] (paths)--(mult);
\draw[arrow] (mult)--(add);
\draw[arrow] (add)--(unk);
\draw[arrow] (unk.west) -| node[pos=0.25,above,font=\small]{да} (yes.north);
\draw[arrow] (unk.east) -| node[pos=0.25,above,font=\small]{нет} (no.north);
\draw[arrow] (yes.south) -- (check.north west);
\draw[arrow] (no.south) -- (check.north east);
\draw[arrow] (check)--(end);
\end{tikzpicture}
\end{center}
\clearpage

\sect{9. Тренировочный блок: 10 задач}
\textcolor{muted}{Решайте последовательно, от базового уровня к более содержательным ситуациям. Во всех задачах указывайте, какие ветви подходят. Для задач 5, 7--10 полезно построить дерево. Ответы приведены отдельно после тренировочного блока.}

\exercise{1}{Два извлечения из коробки}
В учебной лаборатории в непрозрачную коробку положили 3 красных и 7 синих одинаковых жетонов. Ассистент наугад достаёт один жетон, записывает его цвет, возвращает жетон обратно и тщательно перемешивает содержимое. Затем из той же коробки независимо достают второй жетон. Найдите вероятность того, что \textbf{первый жетон окажется красным, а второй — синим}. Объясните, почему условие о возвращении жетона имеет значение.

\exercise{2}{Две шахматные партии}
Два участника школьного шахматного турнира договариваются сыграть две партии. Один из них, игрок А, в первой партии с одинаковой вероятностью получает белые или чёрные фигуры, а во второй обязательно меняет цвет. Из статистики известно, что игрок А выигрывает отдельную партию белыми с вероятностью \(0{,}7\), а чёрными — с вероятностью \(0{,}5\). Считайте исходы партий независимыми при известном цвете фигур. Найдите вероятность того, что игрок А \textbf{выиграет обе партии}.

\exercise{3}{Стрелковая тренировка}
Во время тренировочного упражнения спортсмен делает четыре последовательных выстрела. По данным тренировок вероятность попадания в мишень при первом выстреле равна \(0{,}8\), при втором — \(0{,}9\), при третьем — \(0{,}6\), при четвёртом — \(0{,}7\). Предположите, что результаты независимы. Тренеру нужна вероятность строго определённой последовательности: \textbf{два первых выстрела мимо мишени, а два последних — попадания}. Вычислите эту вероятность, сохранив точность десятичной записи.

\exercise{4}{Четыре последние цифры номера}
Телефонная система случайным образом формирует последние четыре цифры служебного номера. Каждая из четырёх позиций независимо заполняется одной из цифр от \(0\) до \(9\), причём все десять цифр равновероятны. Нуль в начале этой последовательности допустим. Найдите вероятность того, что \textbf{произведение четырёх выбранных цифр будет чётным}. Учитывайте все подходящие случаи, в том числе те, когда чётных цифр несколько.

\exercise{5}{Маршрут к обзорной площадке}
В парке проложены пешеходные дорожки, которые на отдельных развилках расходятся в разные стороны. Турист принимает решения случайно: на каждой развилке выбирает одну из доступных дорожек с одинаковой вероятностью и назад не возвращается. Для попадания на обзорную площадку ему нужно сначала выбрать \textbf{одну определённую} дорожку из трёх, затем \textbf{одну определённую} из двух, а на последней развилке \textbf{одну определённую} из четырёх. Другие выборы приводят к иным выходам из парка. Найдите вероятность, что турист с первого раза попадёт именно на обзорную площадку.

\exercise{6}{Совпадение двух сигналов}
На стенде проверки связи устройство дважды независимо генерирует световой сигнал. Каждый сигнал бывает синим с вероятностью \(0{,}7\) и красным с вероятностью \(0{,}3\). После первого сигнала устройство полностью возвращается в исходное состояние, поэтому цвет второго сигнала не зависит от первого. Инженер считает опыт удачным, когда \textbf{оба сигнала имеют одинаковый цвет}: оба синие или оба красные. Найдите вероятность удачного опыта. Укажите обе подходящие последовательности.

\exercise{7}{Поставка электрических ламп}
Сеть хозяйственных магазинов закупает лампы у двух независимых производителей. Первый завод обеспечивает \(55\%\) всех поступающих ламп, второй — оставшиеся \(45\%\). При выходном контроле первого завода брак составляет \(2\%\), а второго — \(5\%\). Покупатель наугад выбирает одну лампу из общей партии, в которой представлены оба производителя в указанных пропорциях. Найдите вероятность, что выбранная лампа \textbf{исправна}. Сначала отделите вероятность выбора завода от вероятности качества лампы с этого завода.

\exercise{8}{Восстановление доли поставщика}
Один интернет-магазин продаёт зарядные устройства двух производителей. Среди изделий первого производителя неисправные составляют \(1\%\), среди изделий второго — \(4\%\). Доля изделий каждого производителя в продаже заранее неизвестна. По итогам большого числа случайных проверок установлено, что вероятность выбрать неисправное зарядное устройство из всей партии равна \(0{,}028\). Какую \textbf{долю всей партии} поставляет первый производитель? Обозначьте эту долю через \(x\), выразите долю второго через \(x\) и составьте уравнение по двум путям к событию «неисправно».

\exercise{9}{Яйца высшей категории}
Два птицеводческих хозяйства поставляют яйца в общий сортировочный центр. Среди яиц из первого хозяйства \(50\%\) относятся к высшей категории, среди яиц из второго — \(20\%\). Поставки смешиваются, после чего одно яйцо выбирают случайно. По данным центра, вероятность того, что выбранное яйцо относится к высшей категории, равна \(0{,}32\). Определите, \textbf{какая часть всех поступивших яиц привезена первым хозяйством}. Не путайте долю хозяйства в общей поставке с долей яиц высшей категории внутри хозяйства.

\exercise{10}{Автоматическая проверка батареек}
На предприятии после сборки каждая батарейка проходит через автоматический контролёр. По статистике неисправны \(2{,}5\%\) изготовленных батареек. Если батарейка неисправна, система распознаёт дефект и отправляет её в брак с вероятностью \(0{,}96\). Если батарейка исправна, контролёр может ошибочно забраковать её с вероятностью \(0{,}02\). Наугад выбранная батарейка прошла контроль. Найдите вероятность того, что \textbf{контролёр отправил её в брак}, независимо от того, было ли это решение правильным. Выполните расчёт с помощью десятичных дробей либо степеней десяти.

\clearpage
\sect{10. Ответы к тренировочному блоку}
\vspace{5mm}
\renewcommand{\arraystretch}{1.45}
\begin{center}
\begin{tabularx}{0.93\linewidth}{@{}c>{\raggedright\arraybackslash}X>{\raggedleft\arraybackslash}p{34mm}@{}}
\toprule
\textbf{№}&\textbf{Проверяемое умение}&\textbf{Ответ}\\
\midrule
1 & Вероятность одного пути; независимость & \(0{,}21\)\\
2 & Два равновероятных порядка цветов & \(0{,}35\)\\
3 & Фиксированная последовательность событий & \(0{,}0084\)\\
4 & Противоположное событие & \(0{,}9375\)\\
5 & Произведение вероятностей на развилках & \(\dfrac1{24}\)\\
6 & Сумма вероятностей двух несовместных путей & \(0{,}58\)\\
7 & Вероятность исправной лампы из двух поставок & \(0{,}9665\)\\
8 & Уравнение с неизвестной долей поставщика & \(x=0{,}4\)\\
9 & Обратная задача о категориях качества & \(x=0{,}4\)\\
10 & Вероятность отбраковки при ошибках контроля & \(0{,}0435\)\\
\bottomrule
\end{tabularx}
\end{center}
\vspace{9mm}
\subsect{Проверка смысла ответа}
\begin{itemize}
\item Любая найденная вероятность должна находиться в пределах от \(0\) до \(1\).
\item Доля, обозначенная через \(x\), также должна удовлетворять \(0\le x\le1\).
\item Если событие достигается несколькими путями, должны быть учтены \emph{все} подходящие несовместные пути.
\item При получении процента умножьте вероятность на \(100\%\): например, \(0{,}0435=4{,}35\%\).
\item В дереве сумма вероятностей, отходящих от одной вершины, должна быть равна единице.
\end{itemize}
\vfill
\begin{center}\small\textcolor{muted}{Лёвин Артём Александрович · эксклюзивно для Марка Гукина}\end{center}
\end{document}
